\begingroup
\XeTeXlinebreaklocale "en"
\normalfont\fontsize{11}{13}\selectfont
\setlength{\parskip}{1pt plus .3pt}
\thispagestyle{opening}
\pdfbookmark[0]{Author\textquotesingle s Note: English Translation}{origins.en}
\phantomsection\label{origins:en}
\markboth{Open Source Evaluatology}{Author\textquotesingle s Note · English Translation}
\noindent{\sffamily\fontsize{10}{14}\selectfont\color{FrameworkTeal}AUTHOR\textquotesingle S NOTE\quad /\quad ENGLISH TRANSLATION}\par
\vspace{10pt}
\noindent{\sffamily\bfseries\fontsize{24}{29}\selectfont\color{HeadingInk}From Open Source Communities\\to Open Source Evaluatology}\par
\vspace{17pt}
% Paragraph 01
I began using open source software long before I began to understand open source communities. As an undergraduate, I was already using Linux, Apache, and Tomcat in my studies and practical work. What struck me then was their usefulness and openness: when I encountered a problem, I could consult the documentation or read the code. I had not yet thought seriously about who maintained these systems, how strangers managed to work together, or what enabled a project to endure.

% Paragraph 02
The China Open Source Conference in 2017 gave those familiar software names another meaning. For the first time, I came into real contact with the people and communities behind the projects. Some wrote code; others organized collaboration, helped newcomers, or carried out maintenance that received little attention beyond the community. Open source began to take shape for me as a network of human relationships. Looking back, that encounter was an important starting point for my research.

% Paragraph 03
Later, a change in my personal circumstances led me to move to the School of Data Science and Engineering at East China Normal University. The new setting brought a practical question: how should we teach and study open source within a school of data science? An interest had to become a direction worth sustained effort, and observations made in communities had to become questions we could investigate and test. I gradually came to see the public records of open collaboration as a rich, real setting for data science.

% Paragraph 04
Shengyu Zhao's arrival was an important step in this exploration. We had met through the open source community. After leaving Alibaba's Open Source Office, he learned that I was forming an open source research team and came to pursue his doctorate under my supervision. We shared an interest and brought different experiences. Community practice gave our questions a starting point, while academic research required us to formulate them clearly and establish sound evidence. The team took shape gradually.

% Paragraph 05
In 2020, we released the GitHub 2019 Digital Annual Report. Working with approximately 546 million public GitHub event records from 2019, we examined the activities of developers, projects, and organizations, including the place of Chinese participants in global collaboration. This was our first systematic attempt to describe the open source ecosystem through data. The attention and positive response encouraged our young team and showed us that communities and industry had a real need for this kind of observation. In the same year, Shengyu published a blog post titled “Open Source Lights Up the Future.” Its title became our laboratory’s enduring motto.

% Paragraph 06
As the report circulated, we began to consider how to preserve the capabilities developed through that analysis. Data collection and preparation, metric definitions, reproducible computation, and continuing updates were all worth making openly available. We gradually organized the report's processes and methods into an open source project. This was the origin of OpenDigger. Moving from a single report toward data infrastructure that others could help maintain gave our research a basis for continued growth.

% Paragraph 07
In 2021, OpenDigger entered incubation in the Mulan open source community, supported by the China Electronics Standardization Institute. It subsequently began providing open source data services to national departments and institutions. Its applications expanded to support Kaiyuanshe's China Open Source Annual Report, SegmentFault's Hacking Force China rankings, the China Open Source Development Blue Book led by the China Open Source Software Promotion Union, and InfoQ's research on open source development. Each collaboration helped us understand more precisely how the evidence required by one user could differ from that needed by another.

% Paragraph 08
Together with our partners, we gradually developed an ecosystem around OpenDigger. OpenLeaderboard supported browsing and comparison; HyperCRX brought data into the code-hosting interfaces people used every day; and OpenGalaxy gave collaboration networks a spatial form that could be explored visually. The OpenSODA competition, developed with the OpenAtom Foundation, OSGraph with Ant, and OpenTalent with Hubei University brought these capabilities into data innovation, relationship exploration, and talent development. Many new research questions grew out of these concrete uses.

% Paragraph 09
The more our tools were used, the harder it became to avoid the questions behind their metrics. How does the amount of work a person performs relate to its effects? How should a decline in a project's activity be interpreted? Whom can a ranking help, and whom might it overlook? More data alone cannot readily answer these questions. Our activity studies, the OpenRank family of algorithms, and subsequent empirical work moved us forward while repeatedly requiring us to examine what our measures meant and where their limits lay.

% Paragraph 10
During this period of exploration, my exchanges with Professor Jianfeng Zhan deepened. We had first met through the BenchCouncil activities he initiated. In November 2022, we held the first International Forum on Open Source Software Ecosystem Benchmarking, and we have participated closely in the Bench conference each year since. Questions about data, metrics, and applications in open source came into conversation there with research on benchmarks, measurement, and evaluation methods.

% Paragraph 11
Professor Zhan's work on Evaluatology prompted me to reconsider what we had already done. We had often started with a dataset or an application and sought an appropriate method. Evaluatology pressed us to make further questions explicit: what is being evaluated, under what conditions is a comparison made, whose experience of the effects matters, and on what basis is value judged? Questions dispersed across algorithms, platforms, and cases began to acquire a connecting theoretical thread.

% Paragraph 12
At the Bench conference in Guangzhou in December 2024, I delivered a keynote on Open Source Evaluatology, and ECNU received the designation of the first International Research Center for (Open Source) Evaluation Science and Engineering. For us, this marked a further connection between years of practice and the theory of Evaluatology. Open Source Evaluatology acquired a clearer research identity: a domain grounded in Evaluatology that systematically investigates contribution, influence, value, and their evaluation in open collaboration.

% Paragraph 13
I particularly value the two-way relationship. Theory helps us identify assumptions that we had left implicit; practice shows us which concepts need refinement and which methods need revision. Does a metric adequately represent sustained maintenance and coordination? Can the same interpretation be retained when an evaluation is used in different organizations? How might people change their behavior after an evaluation is published? Bringing such questions back to theory creates opportunities for domain research and foundational inquiry to advance together.

% Paragraph 14
Around the time of Bench 2025, the English monograph Evaluatology: The Science of Uncovering the Effects was released; a Chinese edition has appeared more recently. Through Professor Zhan's sustained efforts, I have seen BenchCouncil grow from a venue for conference exchange into an academic ecosystem supported by journals, monographs, and benchmark platforms. A field needs people who pose questions, and it also needs shared settings in which knowledge can accumulate, disagreements can be discussed, and methods can be tested. This development has been an important source of inspiration for us.

% Paragraph 15
Open Source Evaluatology: Theory, Methods, and Empirical Studies took shape in this setting. Writing it required us to revisit years of papers, doctoral dissertations, tools, and applications. It also required an answer to a harder question: why do these achievements belong to a single research field? Across four parts, from theoretical foundations and measurement methods to empirical studies, think-tank applications, and disciplinary development, the book attempts a coherent answer. We want readers to understand how a method works and to judge which questions it is suited to answering.

% Paragraph 16
Developing theory from practice requires restraint. Our cases can provide evidence, but they cannot speak for every open source community. An algorithm's performance on the data available to us does not automatically establish its explanatory power in every setting. We have written this book to make the evidence, methods, and limitations clearer, so that others can examine and extend the work. This is how we hope the foundations of Open Source Evaluatology will become stronger.

% Paragraph 17
My coauthors are Shengyu Zhao and two other doctoral students of mine, Fanyu Han and Xiaoya Xia. Xiaoya has completed her doctorate and joined Ant's Open Source Office, continuing her active engagement in international open source communities. Fanyu is expected to complete his doctorate in 2027. From encounters in the community to shared research and the writing of this book, their growth has been intertwined with that of the team. Being able to share authorship with them is something I deeply value as a teacher.

% Paragraph 18
The four names on the cover cannot represent everyone who has contributed. Students' work on data, code, experiments, and manuscripts; support from colleagues in the school; and the questions, feedback, and opportunities offered by industry partners and community members have all made this sustained effort possible. Many contributions have no name on the cover but remain present in the book. I offer my sincere thanks to everyone who has accompanied us.

% Paragraph 19
Looking back, open source has also helped me answer the question I faced on arriving at a school of data science. Real settings generate needs; applications drive methodological innovation; and innovations return to those settings to be tested and to support further development. The principle that “scenarios drive applications, applications drive innovation, and innovation drives development” has acquired concrete meaning in our work. A report can give rise to a data project, a tool can enter a course or a community, and a question from practice can become a research topic. Teaching, research, and service have found ways to connect.

% Paragraph 20
Through this work, we have provided data and evaluation services to developers, companies, communities, foundations, and government institutions, while helping to educate around thirty young professionals. They have taken different paths, continuing in research or moving into industry and community work. For a university laboratory, producing work that others can use and educating people willing to contribute to shared endeavors are tangible achievements.

% Paragraph 21
Our work has unfolded alongside the development of open source in China. Open source is included in both the Fourteenth and Fifteenth Five-Year Plans, with attention extending from support for communities to the development of an open source ecosystem. Its public significance has become increasingly clear. Open releases by teams such as DeepSeek have also shown how Chinese technology teams can share their work, attract participation from developers worldwide, and build relationships across geographical boundaries. As researchers, we have an opportunity to take part and a responsibility to help make these changes better understood.

% Paragraph 22
On 9 July 2026, at a side event of the AI for Good Global Summit in Geneva, Switzerland, I represented the CAST UN Consultative Committee on Open Source and Innovation (CCOSI) in the release of the 2026 Global Open Source Development Report, Preview Edition. The interest and discussion around the report brought home to me how open source data research could serve a broader international conversation. From analyzing GitHub data for 2019 to preparing a report on global open source development, the scope of our research has expanded. The original wish to understand collaboration has remained.

% Paragraph 23
International dialogue also makes the responsibilities of evaluation more concrete. Conditions of participation differ across regions, and communities leave different kinds of data. An apparently orderly comparison may conceal consequential differences. We hope that scientifically grounded evaluation can help people understand contributions, identify where support is needed, and establish a trustworthy basis for collaboration. Making open work easier to use, and improving the recognition of sustained maintenance and contributions to the public good, are reasons to continue this research.

% Paragraph 24
The development of artificial intelligence makes these questions more pressing. Wider access to technology does not give everyone the same opportunities to participate. Enabling developers from different backgrounds to join collaborative work, sustaining valuable public resources, and extending the benefits of technological progress all require sustained effort. I hope Open Source Evaluatology can contribute a sound way of investigating these questions, giving aspirations for openness and inclusion a basis that can be discussed, tested, and improved.

% Paragraph 25
This book is written in English for an international readership. I nevertheless wanted to record its origins in Chinese and provide an English translation. Chinese is the mother tongue of all four authors. It carries many of the details through which we understand the world, exchange ideas, and express gratitude. Keeping the original is a way of allowing readers to see, alongside the concepts, models, and figures, where this work began, whom we encountered, and why we have continued.

% Paragraph 26
From first using open source software to joining its communities and exploring Open Source Evaluatology with colleagues, what has mattered most has been meeting people and developing shared work from those encounters. I hope this book can make a contribution to that continuing effort. I also hope its readers will carry the work further: helping the value of open collaboration to be better understood, creating more opportunities for those who wish to participate, and sustaining the growth of the open source ecosystem on which we all depend.

\par\vspace{6pt}{\raggedleft\sffamily\fontsize{10}{13}\selectfont Wei Wang\quad Shanghai, November 2026\par}
\clearpage\endgroup
% Restore the original English book line-breaking behavior after bilingual matter.
\XeTeXlinebreaklocale ""
